\section{符号说明}

下表收录全文跨章节沿用的核心记号，一次性出现的中间量在首次出现处就地说明。

\begin{longtable}{clcl}
\toprule
符号 & 含义 & 单位 & 取值或来源 \\
\midrule
\endfirsthead
\toprule
符号 & 含义 & 单位 & 取值或来源 \\
\midrule
\endhead
\bottomrule
\endfoot
$r$ & 到药材中心的径向距离 & m & $[0,R(t)]$ \\
$t$ & 烘干时间 & s & $[0,t_{\mathrm{dry}}]$ \\
$R_0$ & 初始半径 & m & $0.02$ \\
$R(t)$ & 时变外半径 & m & 附件 2，$[0.01198,0.02]$ \\
$L$ & 药材长度 & m & $0.25$ \\
$T(r,t)$ & 药材温度场 & °C & 待求，初值 $28$ \\
$C(r,t)$ & 药材水分浓度（干基含水率） & kg/kg & 待求，初值 $2.55$ \\
$T_{\mathrm{air}}(t)$ & 烘房空气温度 & °C & 附件 1 \\
$W_{\mathrm{air}}(t)$ & 烘房空气水分浓度 & kg/kg & 附件 1 \\
$\rho(C)$ & 药材密度 & kg/m³ & 附录 2/3/4 \\
$c_p(C)$ & 比热容 & J/(kg·K) & 附录 2/3/4 \\
$k(C)$ & 热传导系数 & W/(m·K) & 附录 2/3/4 \\
$D(C,T)$ & 水分浓度扩散系数 & m²/s & 附录 2/3/4 \\
$h$ & 对流换热系数 & W/(m²·K) & $25$，四问共用 \\
$h_m$ & 对流传质系数 & m/s & $8\times10^{-7}$，四问共用 \\
$s$ & 干质量拉格朗日坐标 $(r/R)^2$ & 1 & $[0,1]$ \\
$\rho_s$ & 干骨架密度 & kg/m³ & $\rho_{s0}(R_0/R)^2$ \\
$\bar C$ & 干质量加权平均含水率 & kg/kg & $\int_0^1C\,\mathrm{d}s$ \\
$C_{\mathrm{th}}$ & 烘干阈值 & kg/kg & $0.15$ \\
$t_{\mathrm{dry}}$ & 烘干时长 & h & 待求 \\
\end{longtable}
